% Insertion fragment for the existing zero-geometry chapter.
% Full proofs and exact arithmetic: reports/lerch-global-phase/article/lerch_global_phase.tex
% New labels use globalphase:. No new global macros are defined here.
\section{The complete index-three Lerch phase diagram}
\label{globalphase:section}
The preceding positive-kernel zero bound can be combined with an exterior
inequality to complete the index-three phase diagram. This supplies the
global exclusion left open in Conjecture 12.1 of the companion report
\emph{Lerch Endpoint Singularities and Stieltjes Zero Bifurcations}.
The full details and an executable exact-arithmetic certificate are in
\emph{A Complete Index-Three Lerch Phase Diagram}, stored with this insertion.

Set
\[
 g(x)=\frac{\log^2x(3-\log x)}{x^2},\qquad
 G(a,\rho)=\sum_{m\geq0}\rho^m g(a+m)=F_{3,1}^{\rho}(a).
\]
For $\rho=1$, $G(a,1)=\gamma_3'(a)$. All displayed $G$-values below are
\emph{unnormalized}; they are not the $a^2G$ values in some earlier tables.
For $0<\rho\leq1$ the classical shift equation becomes
\begin{equation}\label{globalphase:shift}
 \rho G(a+1,\rho)=G(a,\rho)-g(a).
\end{equation}
The positive-kernel bound gives at most three positive zeros, counted with
multiplicity. At each fixed $0<a<e^3$, the coefficients in the $\rho$ series
change once from positive to negative, ignoring zero coefficients.
Consequently every interior parameter zero has $G_\rho<0$; positivity at
$\rho=1$ persists for smaller positive parameters.

\begin{lemma}[Global-exclusion inputs]\label{globalphase:inputs}
The following statements hold:
\begin{align*}
 G(a,\rho)&>0 &&\text{for }a\in(0,4/5]\cup[3/2,9/5],\quad0<\rho\leq1,\\
 G_a(a,\rho)+e^{-3}G(a,\rho)&<0 &&\text{for }a\geq e,\quad0\leq\rho\leq1,\\
 G(a,\rho)&>0 &&\text{for }9/5\leq a\leq11/4,\quad0\leq\rho\leq9/10,\\
 G_a(a,\rho)&<-1/40 &&\text{for }9/5\leq a\leq11/4,\quad9/10\leq\rho\leq1.
\end{align*}
There are no left-hand zeros for $0<\rho\leq9/10$.
\end{lemma}
\begin{proof}
The four outward rational enclosures
\begin{align*}
 G(4/5,1)&\in[0.2560298104220831,0.2560298104220832],\\
 G(1,1)&\in[-0.0323050999463418,-0.0323050999463417],\\
 G(3/2,1)&\in[0.0249055835689205,0.0249055835689206],\\
 G(9/5,1)&\in[0.0052639818935007,0.0052639818935008]
\end{align*}
with the three-zero multiplicity bound give the common zero-free regions.
For the exterior inequality, put $y=\log x$ in $g'(x)+e^{-3}g(x)$.
After removal of the positive factor $x^{-3}y$, its sign is that of
\[
 R_3(y)=2y^2-9y+6+e^{y-3}y(3-y),
\]
which is strictly negative for $y\geq1$. One can prove this by writing
$y=3+r$ and using $e^r\geq1+r$ when $r\geq0$, or by writing $y=3-t$
and using $e^{-t}\leq(1+t)^{-1}$ when $0\leq t\leq2$. Summing gives the
second assertion.

For the low-parameter assertion, keep the first positive summand and bound
the negative tail geometrically. The lower bound is
\[
 \frac6{121}-\frac{33}{160}(9/10)^{18}>0.
\]
On $4/5\leq a\leq3/2$, keep the $m=1$ summand instead; the same bound
multiplied by $\rho$ excludes left-hand zeros.

The last inequality is a finite exact certificate, not a sampled numerical
claim. The accompanying verifier partitions the closed $a$-interval into
128 rational cells, encloses 64 derivatives $g'(a+m)$ on each cell, and
bounds the tail by a positive decreasing logarithmic majorant. A nonpositive
upper endpoint is weighted by $(9/10)^m$; a positive endpoint is weighted
by $1$. Thus every bound covers the complete parameter interval. The largest
upper bound, including the infinite tail, is less than
$-0.0256107476350180<-1/40$. Rational logarithm bounds, the anchor
Euler--Maclaurin remainder, every cell, and the independent tests are
specified and reproduced in the full report and its standard-library code.
\end{proof}

\begin{theorem}[Complete global phase diagram]\label{globalphase:theorem}
For every $0<\rho\leq1$ there is exactly one simple outer zero
$a_{\mathrm O}(\rho)\in(9/5,e^3)$, with $G_a<0$ there. There is a unique
$\rho_{\mathrm c}\in(9/10,1)$ at which a double zero
$a_{\mathrm c}\in(4/5,3/2)$ occurs. Below this parameter there is just one
positive zero; at it there are a double zero and the simple outer zero;
above it there are exactly three simple positive zeros. The separate
endpoint $\rho=0$ has the double zero $1$ and the simple zero $e^3$.

At the fold $G_\rho<0$ and $G_{aa}>0$. Above the fold, writing
$a_{\mathrm L}<a_{\mathrm M}<a_{\mathrm O}$, the left and outer branches
strictly decrease, while the middle branch strictly increases, on
$(\rho_{\mathrm c},1)$.
\end{theorem}
\begin{proof}
Existence of an outer zero follows from positivity at $9/5$ and negativity
at $e^3$. For $\rho\leq9/10$, positivity persists to $11/4>e$ and the
exterior integrating factor permits at most one subsequent crossing.
For larger $\rho$, $G$ strictly decreases from $9/5$ to $e$, and
$e^{a/e^3}G$ strictly decreases thereafter. Again only one crossing is
possible. It has negative derivative in either region. Thus no extra
outer pair can occur.

All other zeros must lie in $(4/5,3/2)$. The parameters producing a negative
value somewhere there form an upper interval: at every parameter zero,
$G_\rho<0$. The low-parameter positivity and the negative anchor at $a=1$
give an interior threshold. At the threshold an interior nonnegative
minimum is zero. Its even multiplicity, together with the outer zero and
the bound of three, forces exactly one double zero and $G_{aa}>0$.
A hypothetical zero below the threshold would create negativity by
$G_\rho<0$, a contradiction. Above the threshold a negative value between
positive endpoints gives two simple left zeros; the multiplicity bound
excludes all others. The signs of $-G_\rho/G_a$ give the branch directions.
\end{proof}

\begin{theorem}[Sharp span and integer resonance]\label{globalphase:span}
In the three-zero regime,
\[
 a_{\mathrm O}(\rho)-a_{\mathrm L}(\rho)\leq1.
\]
Equality holds at exactly one parameter $\rho_*$ in
$(\rho_{\mathrm c},1)$, with
$a_{\mathrm L}(\rho_*)=1$ and $a_{\mathrm O}(\rho_*)=2$.
The maximum is nondegenerate. If $D=a_{\mathrm O}-a_{\mathrm L}$ and
$v=a_{\mathrm L}'(\rho_*)$, then
\[
 a_{\mathrm O}'(\rho_*)=v<0,\qquad
 D''(\rho_*)=\frac{6v^2}{G_a(1,\rho_*)}<0.
\]
\end{theorem}
\begin{proof}
The shift equation gives
$G(a_{\mathrm L}+1,\rho)=-g(a_{\mathrm L})/\rho\leq0$.
Since $a_{\mathrm L}+1>9/5>a_{\mathrm M}$, the known sign pattern forces
$a_{\mathrm L}+1\geq a_{\mathrm O}$. Equality is possible exactly when
$a_{\mathrm L}=1$. The unique zero of $G(1,\rho)$ lies in $(0,1)$ because
its first nonzero coefficient is positive and $G(1,1)<0$; the shift then
also gives a zero at $2$. It cannot be the fold, since $g'(1)=0$ would make
both $1$ and $2$ multiple zeros. The span bound shows that $1$ must be the
left rather than middle root. This proves the unique equality case.
Differentiating the shift equation at the simultaneous zeros gives their
common velocity. Finally $g(1+u)=3u^2+O(u^3)$ and the shift equation give
$D''(\rho_*)=6v^2/G_a(1,\rho_*)$ by Taylor expansion at the outer simple root.
\end{proof}

Angle brackets now denote differentiation in the order of the polylogarithm:
$\operatorname{Li}_s^{\langle j\rangle}(\rho)=\partial_s^j
\operatorname{Li}_s(\rho)$. The resonance is the unique interior solution of
\[
 \operatorname{Li}_2^{\langle3\rangle}(\rho_*)
       =-3\operatorname{Li}_2^{\langle2\rangle}(\rho_*).
\]
Put $C=\operatorname{Li}_1^{\langle3\rangle}
+3\operatorname{Li}_1^{\langle2\rangle}$ and
$T=2\operatorname{Li}_3^{\langle3\rangle}
+9\operatorname{Li}_3^{\langle2\rangle}
+6\operatorname{Li}_3^{\langle1\rangle}$. The exact formulas are
\[
 v=\frac{C(\rho_*)}{\rho_*T(\rho_*)},\qquad
 D''(\rho_*)=-\frac{6C(\rho_*)^2}{\rho_*T(\rho_*)^3}.
\]
The full global-phase report gives broader all-index exterior barriers, integer-gap
rigidity, and tail-moment inequalities. A stronger assertion that $D$ is
strictly increasing before $\rho_*$ and strictly decreasing after it remains
an explicitly separate conjecture; the unique global maximum does not
settle that assertion.
